% Transcription — fonds Alexandre Grothendieck, shelfmark 63, batch 2 (pages 21-40).
% Established from the facsimile archives/batches/63/batch-02.pdf (a copy of
% 63.pdf fetched through the project relay, no cover sheet: PDF page k =
% archivists' page 20+k), per the skill transcribe-grothendieck. The folder is
% a hundred pages in five batches; this is the second. Aligned with the
% decisions of batch 1 (transcripts/63/batch-01.fr.tex). Folder 62
% (transcripts/62/) is the notational reference for the group « Courbes
% elliptiques » (62-64).
%
% Pass: Opus 5.5 (claude-opus-5-5), 2026-09-26 — first pass, unchecked against the pages by a human.
%
% Pages skipped: 34, a blank sheet.
%
% Pages given one \note each, not re-set (someone else's text):
% - Pages 21-33: typed pages 7-19 (« - 7 - » ... « - 19 - ») of the
%   typescript « COURBES ELLIPTIQUES : FORMULAIRE d'après J. Tate mis au goût
%   du jour par P. Deligne », whose title leaf is page 15 (batch 1): the end
%   of §2 (after Prop. 2.5), §3 « Cas où 2 et c4 sont inversibles », §4
%   « Cas où 2 est nilpotent et c4 inversible », §5 « Cas où c4 = 0, p = 2 ou
%   p = 3 », §6 « Applications » (Prop. 6.1-6.9) and §7 « Anneau des formes
%   modulaires entières » (Prop. 7.1, 7.2), ending mid-page on typed p. 19.
%   Evidence that it is not his text: the title leaf names Deligne as the
%   author, and the text is the one published as P. Deligne, « Courbes
%   elliptiques : formulaire d'après J. Tate », Modular Functions of One
%   Variable IV, LNM 476 (1975), 53-73 (batch 1's identification; the
%   section titles and proposition numbers here run on without a break from
%   batch 1's pp. 15-20). By the folder rule, one \note per page, with every
%   hand mark given in full. The marks are in pencil: proof-reader's circles
%   (accents, letters, punctuation), prime factorisations (« 12^3 » under
%   1728), and a few worded notes in a fast French hand that looks like his
%   — « section unité ? » (p. 21), j = 0 <=> c4 = 0, j = 12^3 <=> c6 = 0 and
%   « net » over « non ramifié » (p. 26), « O_S » (p. 27), « faux si S non
%   connexe » and « structure ? » (p. 29), « pourquoi ? » and a boxed N.B.
%   on the divisors of c4, c6, Delta on the modular site (p. 31), the
%   smoothness of the divisors of a1 and b2 (p. 32), and five lines of
%   questions (Hasse-Witt, modular forms over an arbitrary field or a DVR,
%   Serre's canonical lifting, the formal group law) on p. 33. The worded
%   notes are set as \marginal{} and attributed to him with the same caution
%   as batch 1's note on p. 20.
% - Pages 35-39: a copy (brown, aged paper) of five handwritten leaves in
%   English, the first unnumbered, then numbered 2 to 5 at the top right, in
%   the same rounded, upright hand as J. Tate's letter copied on pp. 2-13
%   (batch 1) — compared side by side with p. 5 — and nothing like his:
%   Tate's older notation (beta_2 ... beta_8, gamma_4, gamma_6, rho for the
%   scaling factor), the Weierstrass form, the transformation rules, and the
%   « special rational canonical forms » for p ≠ 2,3, p = 3 and p = 2 with
%   their twists and automorphism groups (« G of order 24 »). No publication
%   is identified for these leaves; they are another author's text all the
%   same, and get one French \note per page, not re-set. No mark in a second
%   hand is distinguishable on them.
% - Page 40: the first page of a typescript copy of A. O. L. Atkin,
%   « Congruences for modular forms » (typed header « atkin 1 »), with
%   « (1968) » added in ink by hand and a scrawled monogram top right; one
%   \note. It continues into batch 3.
%
% His pagination: none of his own. The typescript carries typed page
% numbers 7-19; Tate's leaves carry his numbers 2-5 (the first unnumbered);
% Atkin's page is « atkin 1 ». Numbered runs are Deligne's ((2.6),
% (3.1)-(3.5), (4.1)-(4.4), (5.1)-(5.4), (6.1)-(6.9), (7.1)-(7.2)) and
% Tate's (Case 1-3). The only date on these leaves is « (1968) » on p. 40,
% by hand, beside Atkin's title.
%
% What the batch is about: reference material for the folder's formulaire —
% Deligne's scheme-theoretic tables over a base S (normal forms when 2, 3 or
% c4 are invertible, char 2 and 3, the invariant j, Isom(E,E'), automorphism
% groups, the ring of integral modular forms Z[c4,c6,Delta]/(c4^3 - c6^2 -
% 1728 Delta) and its reductions mod 2 and 3), Tate's handwritten canonical
% forms, and the first page of Atkin's paper. His own contribution in this
% batch is limited to the pencil marks on the typescript, given in full;
% the questions on p. 33 (Hasse-Witt interpretation of a1, b2; modular forms
% over any field or a DVR; Serre's canonical lifting by universal formulae;
% explicit formal group law) are the most substantive.
%
% \ill{}: 2. \uncertain{}: 6 (all in the pencil marks, pp. 21, 29, 31, 33).
% Hardest: the notes on p. 29 and the last line on p. 33 (the symbol after
% « Loi de groupe formel sur »).

\documentclass[11pt,a4paper]{article}
% Le préambule commun à toutes les transcriptions du fonds Grothendieck.
%
% Il tient en une idée : **l'incertitude se compose, elle ne se résout pas.**
% Une transcription qui lisse un mot douteux a détruit la seule chose pour
% laquelle on la faisait — un lecteur doit pouvoir savoir, à chaque endroit,
% ce qui était sur le feuillet et ce qui a été ajouté. Les six macros qui
% suivent sont donc l'appareil critique, et non de la décoration.
%
% Les mêmes commandes sont reconnues par `scripts/render.mjs`, qui produit la
% vue de lecture du site. Une macro ajoutée ici doit l'être aussi là-bas,
% faute de quoi le rendu échouera bruyamment — ce qui est voulu : un rendu
% silencieusement faux est pire qu'un rendu absent.

\ProvidesPackage{grothendieck}[2026/08/08 Transcriptions du fonds Grothendieck]

\usepackage{amsmath,amssymb,amsthm}
\usepackage{mathrsfs}
\usepackage{stmaryrd}
% Les diagrammes commutatifs sont partout dans ce fonds : c'est la langue
% même de la géométrie algébrique de Grothendieck, et une transcription qui
% les rendrait en prose ne transcrirait rien.
\usepackage{tikz-cd}
% Français par défaut — c'est la langue des feuillets — mais l'anglais est
% chargé aussi : la traduction et le résumé sont des documents anglais, et la
% typographie française (espaces avant les deux-points) y serait une faute.
% Ces documents basculent par \selectlanguage{english} en tête de corps.
\usepackage[english,french]{babel}
\usepackage{fontspec}
\usepackage{geometry}
\geometry{a4paper,margin=25mm,marginparwidth=28mm}
\usepackage{marginnote}
\usepackage{xcolor}
% ulem, not soul. `soul`'s \st and \ul are built on a character-by-character
% scan that XeLaTeX's font machinery breaks: the document fails with an
% undefined control sequence rather than degrading. ulem does the same two jobs
% and is Unicode-engine safe. `normalem` stops it hijacking \emph, which the
% transcriptions use for Grothendieck's own emphasis.
\usepackage[normalem]{ulem}
\usepackage{hyperref}
\hypersetup{colorlinks=true,linkcolor=black,urlcolor=black,pdfborder={0 0 0}}

% --- Métadonnées du lot -----------------------------------------------------
% Reprises telles quelles par le script de rendu, qui les affiche en tête de
% la vue de lecture. Elles fixent ce que le fichier prétend couvrir : sans
% elles, une transcription flotte sans point d'attache dans un fonds de seize
% mille feuillets.

\newcommand{\folder}[1]{\gdef\grothFolder{#1}}
\newcommand{\batch}[1]{\gdef\grothBatch{#1}}
\newcommand{\leaves}[2]{\gdef\grothFirst{#1}\gdef\grothLast{#2}}
\newcommand{\foldertitle}[1]{\gdef\grothFoldertitle{#1}}
\gdef\grothFolder{?}\gdef\grothBatch{?}\gdef\grothFirst{?}\gdef\grothLast{?}\gdef\grothFoldertitle{}

% --- Le résumé --------------------------------------------------------------
%
% Il ouvre la lecture modernisée, sous le titre « Résumé », et il est composé
% en petit. Ce n'est pas un ornement : c'est le seul passage du document qui
% ne soit pas des mathématiques, et un lecteur doit pouvoir le reconnaître
% avant de l'avoir lu — puis le sauter s'il connaît déjà le sujet, ou s'y
% arrêter s'il ne le connaît pas. Le filet qui le suit marque la reprise du
% corps.

\newenvironment{resume}
  {\small\setlength{\parskip}{0.45em}}
  {\par\medskip\hrule\medskip}

% --- Mots-clés ---------------------------------------------------------------
%
% En anglais, en clôture du résumé de la lecture modernisée : le vocabulaire
% moderne sous lequel chercher ce que les feuillets construisent. Le manifeste
% du site (`npm run manifest`) les extrait pour étiqueter la cote entière —
% cette ligne est la source unique des tags, il n'y a pas de fichier à côté.
\newcommand{\keywords}[1]{\par\smallskip\noindent
  {\footnotesize\textsc{Keywords} --- \textit{#1}}\par}

% --- L'appareil critique ----------------------------------------------------

% Le numéro de feuillet, en marge. C'est l'ancre qui permet de régler un
% désaccord sur une formule en regardant *un* feuillet plutôt qu'une chemise
% de six cents. Le site s'en sert aussi pour tourner les pages du fac-similé
% quand on fait défiler la transcription.
\newcommand{\leaf}[1]{%
  \marginnote{\footnotesize\color{gray!70}\textsf{#1}}%
  \label{leaf:#1}\ignorespaces}

% Illisible : on ne devine pas. Un mot inventé qui se lit comme les autres est
% le pire résultat possible.
\newcommand{\ill}{\textcolor{red!60!black}{[\,\ldots\,]}}

% Lecture douteuse : le mot est donné, et signalé comme incertain.
\newcommand{\uncertain}[1]{\uline{#1}}

% Ajout de l'éditeur — une lettre manquante, un mot sous-entendu.
\newcommand{\add}[1]{\textcolor{blue!50!black}{[#1]}}

% Biffé par Grothendieck lui-même. Ce qu'il a rayé fait partie du document :
% une rature dit souvent où la pensée a changé de direction.
\newcommand{\struck}[1]{\sout{#1}}

% Note du transcripteur, sur l'état matériel du feuillet.
\newcommand{\note}[1]{\par\smallskip{\small\color{gray!60!black}\textit{[#1]}}\par\smallskip}

% Note marginale de Grothendieck, distincte du corps du texte.
\newcommand{\marginal}[1]{\par\smallskip\noindent\hspace{1em}%
  {\small\color{gray!60!black}#1}\par\smallskip}

% --- Titre ------------------------------------------------------------------

\AtBeginDocument{%
  \begin{center}
    {\large\bfseries Fonds Alexandre Grothendieck --- cote n\textsuperscript{o}~\grothFolder}\\[2pt]
    {\itshape\grothFoldertitle}\\[4pt]
    {\small feuillets \grothFirst--\grothLast\ (lot \grothBatch)}\\[2pt]
    {\footnotesize\color{gray!60!black}Transcription établie d'après le fac-similé numérisé
     par l'Université de Montpellier.\\
     Les lectures incertaines sont soulignées, les illisibles notées [\,\ldots\,],
     les ajouts entre crochets.}
  \end{center}
  \vspace{1em}}

\endinput


\folder{63}
\batch{2}
\pages{21}{40}
\dating{[à partir de 1964-vers 1970]}
\watermark{Édition de démonstration}
\foldertitle{Formulaire courbes elliptiques : notes et copies de notes manuscrites (s.d.), tapuscrit (s.d.), copies d'articles (s.d.)}

\begin{document}

\page{21}
\note{Les pages 21 à 33 sont la suite (pages dactylographiées 7 à 19) du
tapuscrit de P.~Deligne « Courbes elliptiques : formulaire, d'après J.~Tate »,
commencé page 15, publié dans \emph{Modular Functions of One Variable IV},
Lecture Notes in Math.~476 (1975), p.~53-73. Texte de Deligne et publié : il
n'est pas recomposé ; chaque page reçoit une note donnant en entier les
marques manuscrites, toutes au crayon. Les quelques notes rédigées sont d'une
écriture rapide qui paraît être la sienne ; l'attribution n'est pas certaine.}
\note{Page 7 du tapuscrit : fin de la Proposition 2.5 (courbe universelle
$Y^2=4X^3-g_2X-g_3$, différentielle invariante $\pi=dX/Y$) ; lorsque 2 et 3
sont inversibles, une forme différentielle invariante rigidifie la courbe ;
(2.6) $c_4^3-c_6^2=1728\,\Delta$ en toute caractéristique ; §3 « Cas où 2 et
$c_4$ sont inversibles » : (3.1) $y^2=x^3+a_2x^2+a_4x+a_6$, les $b_i$, $d_6$,
(3.2) $c_4$, $c_6$ et $\Delta=2^4\delta$. Au crayon : $12^3$ sous le
$1728$ de (2.6) ; en haut à droite, en face de « avec pour courbe universelle
(en coordonnées non homogènes) », une note soulignée :}
\marginal{\uncertain{section} unité ?}

\page{22}
\note{Page 8 du tapuscrit : $\delta$ discriminant du second membre de (3.1) ;
lorsque $c_4$ est inversible, un seul choix de $x$, $y$ donne $d_6=0$ ; (3.3)
$c_6$ et $\Delta$, (3.4) $a_2c_4$, $a_4c_4^2$, $a_6c_4^3$ ; Proposition 3.5 :
au-dessus de $\mathrm{Spec}(\mathbf{Z}[2^{-1}])$, schéma de modules des
courbes de genre un munies d'une forme différentielle invariante inversible
avec $c_4$ inversible, $\mathrm{Spec}(\mathbf{Z}[2^{-1}][a_2,a_4,a_6,
(a_2^2-3a_4)^{-1}]/(a_2a_4-9a_6))$, et sa courbe universelle. Au crayon :
marques de correction (le « s » de « modules » et la fin de « homogènes »
entourés) ; un trait dans la marge gauche en face de « est ».}

\page{23}
\note{Page 9 du tapuscrit : $\pi=dx/2y$ ; §4 « Cas où 2 est nilpotent et
$c_4$ inversible » : (4.1) $y^2+xy=x^3+a_2x^2+a_6$, $\pi=dx/(x+2y)$ ;
(4.2) $r=t=0$, $u=1+2s$ ; (4.3) ; (4.4) modulo 2, $c_4=1$, $c_6=1$,
$\Delta=a_6$ ; §5 « Cas où $c_4=0$, $p=2$ ou $p=3$ », a) $p=3$. Au crayon :
deux marques de correction entourées (la ponctuation après « $p=3$ » du
titre du §5, et le « s » final de « sous »).}

\page{24}
\note{Page 10 du tapuscrit : en caractéristique 3, $c_4=0$ si et seulement si
$a_2^2=0$ ; (5.1) $y^2=x^3+a_4x+a_6$, $c_4=c_6=0$, $\Delta=-a_4^3$ ; (5.2) ;
b) $p=2$ : $c_4=a_1^4$ ; (5.3) $y^2+a_3y=x^3+a_4x+a_6$,
$\pi=dx/a_3=dy/(x^2+a_4)$, $\Delta=a_3^4$ ; (5.4). Au crayon : un trait
entoure le dénominateur $x^2+a_4$ de (5.3) et le rattache, par une flèche, à
la barre de fraction, pour le mettre entre parenthèses ; la ponctuation après
« b) $p=2$ » et la virgule après $\pi$ dans « $(\pi, x, y)$ » sont entourées.}

\page{25}
\note{Page 11 du tapuscrit : §6 « Applications ». Proposition 6.1 : sur un
corps algébriquement clos, $E$ est de type elliptique (resp. multiplicatif,
resp. additif) si et seulement si $\Delta\neq 0$ (resp. $\Delta=0$ et
$c_4\neq 0$, resp. $\Delta=c_4=0$), démonstration a)-d) ; pour $E$ sans fibre
de type additif, définition de la section $j$ de $\mathbf{P}_1/S$. Au
crayon : le « et » entre $\omega^{\otimes 12}$ et $\omega^{\otimes 4}$ est
entouré ; « $P_1$ » est entouré, avec en dessous, entouré, « $P^1$ ».}

\page{26}
\note{Page 12 du tapuscrit : (6.2) $j=c_4^3/\Delta=1728+c_6^2/\Delta$ ;
Proposition 6.3 : pour $E$, $E'$ sans fibre additive, (I) $\mathrm{Isom}(E,E')$
représentable, fini et non ramifié sur $S$, (II) surjectif sur $S$ si
$j_E=j_{E'}$, (III) revêtement étale de rang 2 si de plus $j$ ne prend pas
les valeurs 0 et 1728 ; démonstration ; début de la Proposition 6.4. Au
crayon : sous chacun des trois « Isom » un petit $S$ (lire
$\mathrm{Isom}_S$) ; « non ramifié » entouré, deux fois, avec au-dessus
« net » ; les virgules finales de (I) et (II) entourées et remplacées par
« ; » ; $12^3$ sous le $1728$ de (III). En haut à droite :}
\marginal{$j=0 \Longleftrightarrow c_4=0$}
\marginal{$j=12^3 \Longleftrightarrow c_6=0$}

\page{27}
\note{Page 13 du tapuscrit : Proposition 6.4 : si 2, $c_4$ et $c_6$ sont
inversibles, $\mathrm{Isom}(E,F)$ s'identifie au schéma défini par
$\lambda^2=c_4c_6'/c_6c_4'$ ; démonstration par (6.4.1), (6.4.2) ; début de
la preuve de (III) en général. Au crayon : des parenthèses ajoutées autour de
$j-1728$ dans « $1/j-1728$ » ; le « $S$ » de « sections de $S$ » entouré, et
dans la marge, entouré, $\underline{O}_S$.}

\page{28}
\note{Page 14 du tapuscrit : fin de la preuve de (III) quand 2 est nilpotent
($1/j=a_6+2R(a_2,a_6)$) ; preuve de (II) sur un corps algébriquement clos
quand $j=0$ ou $j=1728$ : a) $p\neq 2,3$, courbes types (6.5)
$y^2=x^3+1$ et (6.6) $y^2=x^3-x$ ; b) $p=3$. Aucune marque manuscrite.}

\page{29}
\note{Page 15 du tapuscrit : (6.7) $y^2=x^3-x$ ; c) $p=2$, (6.8)
$y^2+y=x^3$ ; Proposition 6.9 sur les automorphismes : (I) seuls
l'identité et la symétrie si $c_4$ et $c_6$ sont inversibles ; (II)
$\mathrm{Aut}(E)\to\mathbf{G}_m$ immersion fermée si 2 et 3 sont inversibles ;
(III) image $\mu_6$ (resp.\ $\mu_4$) ; sur un corps, (IV) si $3=c_4=c_6=0$,
produit semi-direct de $\mathbf{Z}/(4)$ par $\mathbf{Z}/(3)$, (V) si
$2=c_4=c_6=0$, groupe d'ordre 24. Au crayon : un trait après
« automorphismes » dans (I) renvoie à une note entourée dans la marge
droite ; le point après « corps » est entouré ; une accolade embrasse (V)
dans la marge, suivie d'une note :}
\marginal{\uncertain{faux si $S$ non connexe}}
\marginal{\uncertain{structure} ?}
\note{la première note est portée en face de (I), la seconde en face de (V).}

\page{30}
\note{Page 16 du tapuscrit : preuve de la Proposition 6.9 : (I) par 6.3 ;
(II) et (III) par l'action sur $\pi$, $c_4=\lambda^4c_4$, $c_6=\lambda^6c_6$ ;
(IV) automorphismes de la forme (6.1) avec $u^4=1$,
$r^3+ra_4+a_6(1-u^2)=0$ et leur loi de composition, $u\in\mu_4$,
$r\in\mathbf{F}_3$ pour (6.7) ; (V) automorphismes de la forme (6.3) avec
$u^3=1$ et deux équations en $s$, $t$. Aucune marque manuscrite.}

\page{31}
\note{Page 17 du tapuscrit : loi de composition des $(u,s,t)$, réduction à
$u\in\mu_3$, $s\in\mathbf{F}_4$, $t^2+t=s^3$ pour (6.8) ; §7 « Anneau des
formes modulaires entières » : définition d'une forme modulaire entière de
poids $n$ ; Proposition 7.1 : l'anneau est engendré sur $\mathbf{Z}$ par
$c_4$, $c_6$, $\Delta$ soumis à la seule relation (7.2)
$c_4^3-c_6^2=1728\,\Delta$ ; début de la démonstration. Au crayon : les
accents circonflexes de « polynôme » ajoutés et entourés (trois fois), le
« d » de « des polynômes » entouré ; dans la marge droite, en face de « Il
reviendrait au même de se limiter aux courbes sans fibres additives »,
soulignée deux fois :}
\marginal{\uncertain{pourquoi} ?}
\note{et, dans un cadre à droite de (7.2) :}
\marginal{N.B. Sur le \uncertain{site modulaire}, les diviseurs de $c_4$,
$c_6$, $\Delta$ sont à multiplicité 1, celui de $\Delta$ est lisse sur
$\mathbf{Z}$ ; ceux de $c_4$, $c_6$ le sont sauf au-dessus des points 2 et
3.}

\page{32}
\note{Page 18 du tapuscrit : Proposition 7.2 : l'anneau des formes
modulaires de caractéristique 2 est $\mathbf{Z}/(2)[a_1,\Delta]$, avec
$c_4=a_1^4$, $c_6=a_1^6$ ; celui de caractéristique 3 est
$\mathbf{Z}/(3)[b_2,\Delta]$, avec $c_4=b_2^2$, $c_6=-b_2^3$ ;
démonstration (formes $a_1^{-n}P(a_1,a_6)$, $a_1^{-m}Q(a_1,\Delta)$). Au
crayon : le « 1 » de « (4.1) » et l'accent de « polynôme » entourés ; sous
chacun des deux énoncés, une flèche partant de $a_1$ (resp.\ $b_2$) et une
ligne :}
\marginal{le diviseur de $a_1$ est lisse sur $\mathbf{F}_2$}
\marginal{le diviseur de $b_2$ est lisse sur $\mathbf{F}_2$}
\note{sic : « $\mathbf{F}_2$ » aussi dans la seconde ligne, qui porte sur la
caractéristique 3.}

\page{33}
\note{Page 19 du tapuscrit, dernière du texte dactylographié ici : fin de la
démonstration de 7.2 (une forme de caractéristique 3 s'écrit
$b_2^{-n}P(b_2,\Delta)$ et est définie pour $j=0$ si et seulement si c'est un
polynôme en $b_2$ et $\Delta$). L'accent de « polynôme » est entouré. Le bas
de la page porte, au crayon, cinq lignes de sa main vraisemblablement, les
deux dernières séparées des trois premières par un trait et réunies par une
accolade :}
\marginal{Interprétation de $a_1$, $b_2$ comme forme de Hasse-Witt ?}
\marginal{Formes modulaires sur un corps quelconque ? Sur un ann.\ de val.\
discrète\ldots}
\marginal{Hasse-Witt sur un corps de car.\ $p>0$ quelconque ?}
\marginal{\uncertain{Relèvement canonique} de Serre par formules universelles}
\marginal{Loi de groupe formel sur \ill{} : formules explicites\ldots}
\note{après « sur », un symbole que nous ne lisons pas, portant un indice qui
ressemble à $M/S$ et une surcharge au-dessus.}

\page{35}
\note{Les pages 35 à 39 sont la copie de cinq feuillets manuscrits en anglais
(le premier sans numéro, les suivants numérotés 2 à 5 en haut à droite), de
la même main que la lettre de J.~Tate copiée pages 2 à 13, et non de la
sienne ; notation plus ancienne ($\beta_2,\dots,\beta_8$, $\gamma_4$,
$\gamma_6$, facteur $\rho$). Texte d'un autre auteur : il n'est pas
recomposé, chaque page reçoit une note. Aucune marque d'une seconde main
n'y est discernable.}
\note{Feuillet 1 : l'équation encadrée
$y^2+a_1xy+a_3y=x^3+a_2x^2+a_4x+a_6$, les $\beta_i$, $\gamma_4$, $\gamma_6$,
$\Delta$, $1728\,\Delta=\gamma_4^3-\gamma_6^2$ ; les transformations
birationnelles $x=\rho^2x'+r$, $y=\rho^3y'+s\rho^2x'+t$ et la transformation
des coefficients, des $\beta_i$ et des $\gamma_i$ ; $\rho^{12}\Delta'=\Delta$,
$\omega'=\rho\,\omega$ ; dans des cadres, $j=\gamma_4^3/\Delta$,
$12g_2=\gamma_4$, $216g_3=\gamma_6$, $(j-1728)\Delta=\gamma_6^2$, et la
courbe « générique » de module $j$, $y^2+xy=x^3-\frac{36}{j-1728}x
-\frac{1}{j-1728}$.}

\page{36}
\note{Feuillet 2 : $F$, $F_x$, $F_y$, la différentielle de première espèce
$\omega$ et sa transformation ; « Special rational canonical form » pour
$p\neq 2,3$, $y^2=x^3+a_4x+a_6$ : $j$ comme invariant absolu, et trois cas
($j\neq 0,1728$ ; $j=1728$ ; $j=0$), chacun avec le critère d'équivalence
de deux courbes de même $j$, l'extension $k'$ qui les rend isomorphes, les
automorphismes et la courbe canonique ($y^2=x^3-x$, $y^2=x^3-1$).}

\page{37}
\note{Feuillet 3 : $p=3$, $y^2=x^3+a_2x^2+a_4x+a_6$, $\Delta$, $j=a_2^6/\Delta$ ;
cas 1, $j\neq 0$, forme $y^2=x^3+a_2x^2+a_6$ ; cas 2, $j=0$, forme
$y^2=x^3+a_4x+a_6$, $k'/k$ séparable de degré $\leq 12$, groupe
d'automorphismes d'ordre 12, courbe canonique $y^2=x^3-x$ ; une ligne
biffée.}

\page{38}
\note{Feuillet 4 : $p=2$ ; cas 1, $a_1\neq 0$, forme
$y^2+xy=x^3+a_2x^2+a_6$, $j=1/a_6$, courbe canonique $y^2+xy=x^3+a_6$ ;
cas 2, $a_1=0$, forme $y^2+a_3y=x^3+a_4x+a_6$, $j=0$, $k'/k$ séparable de
degré $\leq 24$, courbe canonique $y^2-y=x^3$, automorphismes
$\rho^3=1$, $s^4+s=0$, $t^2+t+s^3+s^2=0$, « G of order 24 ».}

\page{39}
\note{Feuillet 5 : « Summarizing » : pour $j\neq 0,1728$, en toute
caractéristique, autant de courbes d'invariant $j$ sur $k$ que d'extensions
quadratiques séparables ; dans tous les cas, deux courbes de même $j$ sont
birationnellement équivalentes sur une extension séparable de degré
divisant 24, et le groupe d'automorphismes est un sous-groupe du groupe
d'ordre 24.}

\page{40}
\note{Première page d'une copie dactylographiée de l'article de
A.~O.~L.~Atkin, « Congruences for modular forms » (en-tête dactylographié
« atkin 1 » ; l'auteur y est « Research Fellow of the Atlas Computer
Laboratory ») : introduction sur l'usage de l'ordinateur Atlas pour les
congruences des formes modulaires. Texte d'un autre auteur, non recomposé.
À la main, à l'encre, « (1968) » entouré à droite du titre ; en haut à
droite, un paraphe \ill{}. La copie se poursuit dans le lot suivant.}

\end{document}
